% ============================================================================
%  Langue as a Lattice
%  Closing the Circle from Statistical to Structural Linguistics
%  through Content-Addressed Algebras
% ----------------------------------------------------------------------------
%  Build:  pdflatex langue_as_a_lattice.tex   (run twice)
%  notation follows main.tex / companion.tex / tecs_state_machine_trajectory.tex
%  / attention_lattice_traversal.tex
% ============================================================================
\documentclass[11pt,a4paper]{article}

\usepackage[utf8]{inputenc}
\usepackage[T1]{fontenc}
\usepackage{amsmath,amssymb,amsthm}
\usepackage{amscd}
\usepackage{booktabs}
\usepackage{array}
\usepackage{hyperref}
\usepackage[margin=1in]{geometry}
\usepackage{enumitem}
\usepackage{xcolor}
\definecolor{hlblue}{RGB}{0,80,160}
\hypersetup{colorlinks=true,linkcolor=hlblue,urlcolor=hlblue,citecolor=hlblue}

% ---- theorems (aligned with the PUB series) --------------------------------
\newtheorem{principle}{Principle}[section]
\newtheorem{theorem}[principle]{Theorem}
\newtheorem{definition}[principle]{Definition}
\newtheorem{proposition}[principle]{Proposition}
\newtheorem{corollary}[principle]{Corollary}
\newtheorem{remark}[principle]{Remark}
\theoremstyle{remark}
\newtheorem{observation}[principle]{Observation}

% ---- shorthands (aligned with the PUB series) ------------------------------
\newcommand{\Plane}{\ensuremath{\mathbb{B}}}            % bit plane
\newcommand{\Lat}{\ensuremath{\mathcal{B}}}              % HLLSet (Boolean) lattice
\newcommand{\Tok}{\ensuremath{\mathcal{T}}}              % token space
\newcommand{\HLL}{\ensuremath{H}}                        % an HLLSet
\newcommand{\xor}{\ensuremath{\Delta}}                   % ring addition
\newcommand{\BSS}{\ensuremath{\mathsf{BSS}}}             % BSS measure
\newcommand{\res}{\ensuremath{\mathrm{res}}}             % residual
\newcommand{\spanof}[1]{\ensuremath{\langle #1\rangle}}  % GF(2) span
\newcommand{\supp}{\ensuremath{\mathrm{supp}}}           % support
\newcommand{\Attn}{\ensuremath{\mathrm{Attn}}}           % attention operator
\newcommand{\Ctx}{\ensuremath{\mathcal{C}}}              % token context
\newcommand{\Ing}{\ensuremath{\mu}}                      % ingestion morphism
\newcommand{\Mat}{\ensuremath{\rho}}                     % materialization morphism
\newcommand{\Lang}{\ensuremath{\mathcal{L}}}             % langue (the lattice system)
\newcommand{\Par}{\ensuremath{\mathcal{P}}}              % parole (usage)
\newcommand{\Sign}{\ensuremath{\mathcal{S}}}             % the sign relation
\newcommand{\Cov}{\ensuremath{\mathrm{cov}}}             % coverage
\newcommand{\Disp}{\ensuremath{\mathcal{D}}}             % the dispatcher

% ---- title -----------------------------------------------------------------
\title{\textbf{Langue as a Lattice}\\[4pt]
\large Closing the Circle from Statistical to Structural Linguistics\\
through Content-Addressed Algebras\\[6pt]
\normalsize Fifth paper in the \textit{Emerging World Models} series}

\author{Alex Mylnikov \and Alexandr Mylnikov Jr. \and Deependra Kumar}
\date{\today}

\begin{document}
\maketitle

\begin{abstract}
Modern language models practise \emph{statistical} linguistics: attention
matrices and distributional semantics describe language as usage---what
Saussure called \emph{parole}. Structural linguistics, from Saussure onward,
claims that behind usage stands a system of pure differences---\emph{langue}---
but that system has resisted direct computation: it has been a philosophical
object, not an executable one. This paper shows that a content-addressed
\emph{morphism}---the IICA map from tokens to algebraic structure---closes the
gap; in the reference implementation this morphism is a hash, chosen only
because it distributes its values evenly over the plane. The map
$\Ing : \Tok \to \Plane$ compiles any token language
into a fixed bit plane, where the HLLSet Boolean algebra and its GF(2) ring
form an \emph{executable system of differences}: novelty and retention are the
syntagmatic step, and the directed soft key $\BSS_w(X)=|X\cap B_w|/|B_w|$
reads succession against a reference state---the compositional tool of the
lattice; the ring basis is the set of distinctive oppositions; ingestion and
materialization are the
two faces of the sign, governed by the coverage invariant; and idempotence,
immutability, and content addressability (IICA) are the social contract that
fixes the otherwise arbitrary sign. The resulting construction---the EWM
lattice plus a learned dispatcher that selects which host language actualizes
the shared langue---closes a circle: from statistical linguistics through a
content-addressed morphism into algebraic structures, and back to a
computational form of fundamental structural linguistics. Along the way the
classical hierarchy of linguistic description---phonology, morphology, syntax,
semantics, pragmatics, discourse---is re-founded level by level in the
lattice and the ring, not as an analogy but as a computation. The dispatcher
is the executive side of language made trainable: langue does not speak,
parole does, and the Emerging LLM learns \emph{who speaks}.
\end{abstract}

\noindent\textbf{Keywords:} structural linguistics; Saussure; langue; parole;
HLLSet; Boolean lattice; GF(2) ring; content addressing; IICA; attention;
dispatcher; materialization; sign.

% ============================================================================
\section{Introduction}
\label{sec:intro}
% ============================================================================

Two descriptions of language have never quite met. The statistical description
says that meaning is distribution: a word is known by the company it keeps
\cite{harris1954,firth1957}, and a modern attention matrix is the strongest
realization of that slogan---at every layer, a token's context is a weighted
sum of the tokens that co-occur with it \cite{vaswani2017}. The structural
description says the opposite: a unit's value is not its use but its
\emph{position in a system of differences}; in language ``there are only
differences without positive terms'' \cite{saussure1916}. The first view is
computable but shallow---it gives usage, not system. The second view is deep
but has seemed non-computable---it gives system, not a machine that runs it.

This paper proposes a bridge, and the bridge is an \emph{IICA morphism}: a
deterministic, immutable, content-addressed map from tokens into algebraic
structure. A hash function is only one possible implementation of that
morphism---the one the reference system uses, for the single reason that it
spreads its values evenly over the plane.

The Emerging World Models (EWM) programme \cite{ewm2026,companion2026,state2026}
already contains the necessary pieces, built for other reasons. Its state is a
content-addressed set---an \textbf{HLLSet}, a bit-vector sketch of a token
collection over a fixed plane $\Plane$ of $32{,}768$ bits---and its morphisms
are the Boolean operations $\cup, \cap, \xor$ and the GF(2) ring they generate.
Tokens are mapped to bits by a morphism $\Ing$ (a hash in the reference
implementation); bits are mapped back to tokens by
\textbf{materialization} $\Mat$ \cite{attention2026}. The three properties that
make this mapping well-defined---idempotence, immutability, content
addressability (IICA) \cite{ewm2026}---are exactly the properties a sign
system needs in order to be \emph{arbitrary yet fixed}. And the newest piece of
the programme, the Emerging LLM dispatcher, is precisely the executive side of
language: a learned policy that reads the lattice (langue) and chooses which
host language (parole) should actualize it.

The paper's thesis is one sentence: \textbf{the HLLSet lattice is a computable
langue}. We develop the claim in five steps. Section~\ref{sec:parole} states
the Saussurean distinction in the language of modern models and identifies
attention as parole-level statistics. Section~\ref{sec:morphism} shows that
the morphism step $\Ing$ is the computable reduction from usage to structure.
Section~\ref{sec:langue} exhibits the lattice as a system of differences:
D/R/N are syntagmatic, the directed soft key BSS reads succession (and, as a
profile over a frame, the paradigmatic axis), and the ring basis is the set of
distinctive oppositions. Section~\ref{sec:sign} identifies
ingestion/materialization with the sign relation and IICA with the social
contract that fixes the arbitrary sign. Section~\ref{sec:dispatch} interprets
the dispatcher as actualization---the selection of parole---and closes the
circle. Section~\ref{sec:relations} situates the result against Saussure,
Hjelmslev, distributionalism, and the categorical reading of the bridge.
Section~\ref{sec:conclusion} summarizes.

% ============================================================================
\section{Parole: attention as statistical linguistics}
\label{sec:parole}
% ============================================================================

Saussure's distinction is between \emph{langue}, the social system of a
language---its units, rules, and oppositions, existing independently of any
speaker---and \emph{parole}, the individual acts of speech that actualize the
system. Parole is heterogeneous, statistical, and executive; langue is
homogeneous, structural, and social. A speaker never \emph{sees} langue; they
see parole, and from parole they acquire langue.

Modern language models are the strongest parole-machines ever built. A
transformer's attention matrix
\[
  \Attn(X) \;=\; \mathrm{softmax}\!\left(\frac{QK^\top}{\sqrt{d}}\right)V
\]
is a differentiable co-occurrence table: for each query token it records how
much of every other token's value is \emph{in use} at this step. The weights
are frequencies, the context is a mixture, and the entire enterprise---masked
language modeling, next-token prediction, retrieval, embeddings---is the
distributional programme of \cite{harris1954,firth1957} carried to its limit.
Attention describes usage with unprecedented fidelity, but it remains a
description of \emph{parole}: a model knows that two tokens co-occur; it does
not know, as a matter of structure, \emph{what their difference consists in}.

The missing object is langue. What would it mean to compute it? Three
requirements follow from Saussure's text. (i) \emph{Differences, not positives:}
a unit is defined by its relations to other units, never by its own content.
(ii) \emph{Arbitrary but fixed:} the link between a signifier and what it
signifies is unmotivated, yet the system fixes it by convention.
(iii) \emph{Social and shared:} the system is the same for all speakers; it is
what makes their divergent utterances commensurable. The next three sections
show that the EWM stack satisfies (i) by construction, (ii) by IICA, and
(iii) by the bit plane being host-agnostic.

% ============================================================================
\section{The morphism step: from usage to structure}
\label{sec:morphism}
% ============================================================================

The first move is the least glamorous and the most decisive. Let $\Tok$ be the
universe of tokens---across \emph{all} host languages, so that the same token
string may be emitted by different systems with different vocabularies. Define
the plane
\[
  \Plane \;=\; \{0,\dots,32767\},
\]
and the ingestion morphism
\[
  \Ing : \mathcal{P}_{\mathrm{fin}}(\Tok) \;\longrightarrow\; \mathcal{P}(\Plane),
  \qquad
  \Ing(X) \;=\; \bigcup_{t \in X} \Ing(\{t\})
\]
where each token maps to a fixed subset of $\Plane$ under a deterministic
morphism $\Ing$. The morphism is what matters; the hash is only its reference
implementation, chosen for a single reason---it distributes the function values
evenly over the plane.

\begin{definition}[IICA morphism]
A map $\Ing : \Tok \to \mathcal{P}(\Plane)$ is an \textbf{IICA morphism} if it
is
(i) \emph{idempotent}, read as \emph{determinism}: the same token always maps
to the same subset, $\Ing(t) = \Ing(t)$;
(ii) \emph{immutable}: once fixed, the image is never rewritten; and
(iii) \emph{content-addressed}: the image is named by its content alone.
Any deterministic map satisfies (i). Among all such maps, a hash is preferred
in practice only because it spreads the images evenly across the plane, which
keeps the resulting bit-vector compact and unbiased.
\end{definition}

The image of a token collection under $\Ing$ is an \textbf{HLLSet}: a compact
bit-vector in which repetition has already collapsed (a token present twice
contributes once), and union is the only accumulation.

\begin{proposition}[Ingestion is join-preserving]
For finite token sets $X, Y \subseteq \Tok$,
\[
  \Ing(X \cup Y) \;=\; \Ing(X) \cup \Ing(Y).
\]
Consequently the bag-of-tokens statistics of parole---co-occurrence, overlap,
containment---all reduce to the single operation of union on the plane, and the
order in which tokens arrive is forgotten.
\end{proposition}

The proposition is the entire content of the ``morphism step''. It is what
statistical linguistics lacks: a canonical quotient that throws away frequency
and sequence, and keeps exactly what is left when one asks \emph{which things
occur together as a set}. Idempotence is the first IICA property and it is
already at work here: because the morphism is deterministic, ingesting the same
content twice is a no-op, so the system's memory is a lattice of \emph{types},
not a pile of \emph{tokens}.

\begin{remark}[Why a morphism, not an embedding]
An embedding maps a token to a dense vector; similarity then lives in a metric
space and must be learned. A morphism into $\mathcal{P}(\Plane)$ maps a token
to a \emph{subset of a shared finite set}; similarity then lives in the Boolean
algebra of that set and is \emph{exact}. The first keeps language in the
statistical world; the second moves it into the structural world. \emph{Which}
deterministic morphism is used is secondary: any IICA morphism works, and the
reference system chooses a hash purely for its even distribution of values.
This is the crossing point of the circle.
\end{remark}

% ============================================================================
\section{Langue as a lattice of differences}
\label{sec:langue}
% ============================================================================

On the plane, structure is no longer metaphorical. The set $\Lat =
\mathcal{P}(\Plane)$ is a Boolean lattice under $\cup$ and $\cap$, and a
Boolean ring under $\xor$ and $\cap$ \cite{birkhoff1940}. Every token
collection has become a node of $\Lat$; every relation between collections has
become a relation between nodes. We now read Saussure's system of differences
off the lattice.

\subsection{Syntagmatic relations: D/R/N}
The syntagmatic axis is the axis of combination---what is present together, in
sequence. In the lattice it is read as the stepwise decomposition of a context
$\Ctx$ against a document $D$:
\[
  N \;=\; |D \setminus \Ctx|, \qquad
  R \;=\; |D \cap \Ctx|, \qquad
  \Cov \;=\; R / |D| .
\]
Novelty $N$ is what the context cannot account for; retention $R$ is what it
confirms; coverage is their ratio. These are pure differences: no token is
valued positively, only as it stands in relation to the context.

\subsection{Directed containment: the soft key BSS}
The second reading is deliberately asymmetric. Against a reference state $B_w$
(a basis vector, a host codebook, or a previous context), a state $X$ is read
as
\[
  \BSS_w(X) \;=\; \frac{|X \cap B_w|}{|B_w|},
\]
the fraction of the reference that $X$ covers. The direction matters: it
measures evolution from one state to another---how much of the earlier state
the later state retains, covers, and composes. The EWM programme chose this
asymmetric form as its primary measurement because the object of study is
succession; the symmetric cosine is retained for the cases that need it.

Linguistics is likewise about sequences: a text is one state after another.
The directed BSS is therefore the compositional/generative tool of the
lattice---it reads a new state against an old one and says how much of the old
is carried forward, which is exactly the syntagmatic succession Saussure
placed on the axis of combination. Read as a profile over a whole frame, the
same soft key is also the paradigmatic reading: each coordinate measures the
query against one opposition.

\begin{remark}[Symmetric readings when needed]
The directed BSS is not the only measurement on the plane. When the question
is static resemblance rather than succession---two hosts, two contexts, two
readings of one window, with no direction between them---the symmetric
cosine over the same supports,
\[
  \mathrm{cos}(X, Y) \;=\; \frac{|X \cap Y|}{\sqrt{|X|\,|Y|}},
\]
is available and exact for the same reason: both measurements live in the
Boolean algebra. They are not rivals. BSS answers \emph{how much of the
reference does the new state carry forward}---composition and generation; the
cosine answers \emph{how alike are these two states}---comparison and
classification. The grammar uses the first; the lexicon, and any comparison of
hosts or contexts on an equal footing, uses the second.
\end{remark}

\subsection{Distinctive oppositions: the ring basis}
The deepest structural object is the basis. Over any window of originals, the
EWM apparatus maintains an Reduced Row Echelon Form (RREF) GF(2) basis $F = (B_1,\dots,B_k)$ whose span is
exactly the set of originals \cite{state2026}. A frame is a system of
distinctive oppositions in the phonological sense: each basis vector is a
minimal unit that cannot be assembled from the others, and a node's hard key
$c \in \mathrm{GF}(2)^k$ is its unique decomposition
\[
  X \;=\; \bigoplus_{i : c_i = 1} B_i,
\]
defined exactly when $X$ lies in the span. The soft key
$w_i = |X \cap B_i| / |B_i|$---the directed BSS of Section~4.2, one coordinate
per opposition---is the relaxed reading of any node against the frame, and the
spill
$s = |X \setminus \mathrm{cover}(F)|$ is the part the system cannot yet see.
This is Saussure's ``differences without positive terms'' as an executable
datatype: every unit is an address in the span of the basis, and nothing else.

\begin{observation}[The meta-language]
The HLLSet algebra is not another token language; it is a \emph{meta-language}.
It describes descriptions. Any host language---any vocabulary, any tokenizer,
any attention geometry---compiles into the same plane, and all comparisons
happen there. Host-agnosticism is therefore not an engineering preference but
the structuralist condition (iii) of Section~\ref{sec:parole}: the system is
shared because it is the only level where different paroles are commensurable.
\end{observation}

\subsection{The classical taxonomy, recomputed}
Classical linguistics organises language into seven levels. Two of them do not
survive the move from acoustic speech to text---phonetics and phonology are
absent from LLMs except in special cases (speech and audio models). The other
five survive, but only when they are \emph{recomputed}: the lattice and the
ring do not mimic the classical levels, they re-found them. The dictionary is
as follows.

\begin{center}
\begin{tabular}{@{}>{\raggedright\arraybackslash}p{2.0cm}
                    >{\raggedright\arraybackslash}p{2.4cm}
                    >{\raggedright\arraybackslash}p{10.4cm}@{}}
\toprule
Level & In LLMs & In the HLLSet lattice / GF(2) ring \\
\midrule
Phonetics & N/A (except speech) & the raw token event---the substance \\
Phonology & N/A (except special cases) & basis frames: distinctive oppositions, hard key \\
Morphology & subword tokens & union/join: composition and derivation; idempotence collapses reduplication \\
Syntax & attention over positions & joins and meets; the ring sum; the span as the set of well-formed constructions \\
Semantics & embeddings & directed BSS (coverage of a reference), containment, the hard key as structural address \\
Pragmatics & instruction tuning & the dispatcher: who speaks, fallback, confidence---actualization \\
Discourse & the context window & the commit DAG, $H(t)$, windows, and $D/R/N$ over time \\
\bottomrule
\end{tabular}
\end{center}

\emph{Phonetics} is the physical layer, and in a text-bound system its only
trace is the raw token event itself---the substance that Hjelmslev would place
on the expression plane. \emph{Phonology} is where the re-foundation begins:
the basis frame is a distinctive-feature system. Each basis vector is a minimal
opposition; the hard key is the phonemic address of a node; the soft key is its
allophonic variation against every opposition; and the spill is the part of an
utterance that falls outside the system's inventory. \emph{Morphology} becomes
set composition: derivation is join, compounding is union, and reduplication is
erased by idempotence---repeating a morpheme contributes nothing twice, exactly
as it should in a system of types. \emph{Syntax} becomes the algebra of
combination without linear order: constituents are joins and meets, the ring
sum is phrase-level difference, and the span of the basis is the set of
well-formed constructions over a window. Linear order is the one thing the
lattice forgets, and it is recovered not by the algebra but by materialization,
which restores order at the door. \emph{Semantics} is the system of values
itself: the directed BSS says how much of a reference meaning a new state
covers, containment and coverage say how much of a meaning a context
determines, and the hard key is the structural address
of a meaning---Saussure's value made literal. \emph{Pragmatics} is the
dispatcher: context-dependent selection of who speaks, confidence as the
speaker's commitment, fallback as repair. \emph{Discourse analysis} is the
trajectory: cohesion is retention $R$, progression is novelty $N$, and the
commit DAG with its windows and crossings is the memory of a conversation.

The mapping is deliberately not exact, and that is the point. The classical
hierarchy was built for spoken language and inherited by text; the lattice and
ring re-found it for computational structural linguistics, where every level is
a computation over one shared plane rather than a separate descriptive
practice.

% ============================================================================
\section{The sign and its contract: ingest, materialize, IICA}
\label{sec:sign}
% ============================================================================

Saussure's sign unites a signifier and a signified; the link is arbitrary, and
the system fixes it by convention. The EWM bridge is the same object, made
precise. The sign relation is the pair
\[
  \Sign \;=\; (\Ing, \Mat),
  \qquad
  \Ing : \Ctx \to \Lat, \qquad \Mat : \Lat \to \Ctx,
\]
where materialization restores, from a lattice node, an ordered token context
through the door \cite{attention2026}. The round-trip is not automatic; it is
guarded:

\begin{proposition}[The round-trip law]
For a context $\Ctx$ and its node $\HLL = \Ing(\Ctx)$,
\[
  \Mat(\Ing(\Ctx)) \;=\; \Ctx
  \qquad\text{up to order,}\qquad
  \text{iff}\qquad \Cov(\Ctx) = 1 .
\]
Coverage $1.0$ is the condition under which the sign relation is faithful:
what was compressed is exactly what can be restored. Below it, the spill and
the residual measure the failure, and the system knows the failure's size.
\end{proposition}

The three IICA properties are then read as the \emph{social contract of the
sign}. \textbf{Idempotence}: repeating a sign changes nothing---the system's
value of a content is stable under repetition. \textbf{Immutability}: the
system never rewrites its own history, so the sign's past values remain
referable. \textbf{Content addressability}: the same content always bears the
same name, in every host language, so the arbitrary link between signifier and
signified is fixed by the system rather than by any speaker. Arbitrary but
fixed: this is exactly Saussure's formulation of the linguistic sign, stated
as a property of morphisms.

\begin{remark}[The categorical reading]
In categorical terms the bridge is an adjunction in embryo: ingestion is a
free/compressive functor from the category of token contexts to the Boolean
lattice, materialization is its right adjoint restricted to the image, and
IICA is the naturality that makes the correspondence independent of the host
language. We sketch this in Section~\ref{sec:relations}; a full treatment is
left to a companion note.
\end{remark}

% ============================================================================
\section{Actualization: the dispatcher closes the circle}
\label{sec:dispatch}
% ============================================================================

Langue never speaks. Parole does. The remaining question is therefore
Saussure's executive question: given a state of the system, \emph{which}
parole should actualize it? In the EWM line this question already has a
computational form---it is the router problem solved, in their own ways, by the
typed-decision models Jev and Laya, and by the structural router of the
programme's notebooks \cite{state2026}.

The Emerging LLM is a different answer to the same question. It is a
\textbf{dispatcher}
\[
  \Disp : \Lat \times \mathrm{State} \;\longrightarrow\; \Delta(\mathrm{Hosts}),
\]
a learned policy that reads the meta-language state---BSS, D/R/N, ring basis,
windows---and produces a distribution over host languages. It does not speak
any host language natively; it speaks the meta-language and selects a native
speaker.

\begin{remark}[Recognition without production]
This is a recognizably human ability. A person can identify a language they do
not speak---French, Italian, Japanese---from its sound and rhythm alone, long
before they can produce a single grammatical sentence in it. The dispatcher
exhibits the same asymmetry, structurally: it recognises host languages from
their traces in the lattice (their codebook projections, their coverage) while
producing none of them itself. Recognition without production, in both cases,
is possible because recognition operates on the system of differences, not on
the ability to actualize it.
\end{remark}

\emph{The circle closes}. Parole (a host's attention and tokens) is
compiled by $\Ing$ into langue (the lattice); langue is read back through
$\Mat$ into token context; the dispatcher reads langue and selects the host;
the host generates, and its output re-enters as new parole. Statistical
linguistics $\to$ content-addressed morphism $\to$ algebraic structure $\to$
the executive side of language $\to$ statistical linguistics again, one level
up.

\begin{observation}[Native labels]
The lattice labels its own decisions. The per-host codebook projections of a
materialized context---the dense association of token/encoding sets to
HLLSet bits---form a labeling system independent of any host's own labels,
emergent with the lattice, and traceable through the HLLSet Ring basis generations. The
dispatcher is therefore trainable against langue itself, with external router
decisions serving only as validation.
\end{observation}

% ============================================================================
\section{Relations}
\label{sec:relations}
% ============================================================================

\paragraph{Saussure.} The mapping is direct: langue is $\Lat$; parole is the
stream of host outputs; the sign is $(\Ing, \Mat)$; the arbitrariness of the
sign is the freedom of the morphism (any deterministic IICA map would do), and
its fixedness is IICA; value through difference is D/R/N, BSS, and the basis
hard key.

\paragraph{Hjelmslev.} The glossematic distinction between expression and
content, and the insistence that both sides have their own form and substance,
is realized by the symmetry of the bridge: token context is the expression
plane, the HLLSet is the content plane, and the morphisms $\Ing, \Mat$ are the
two directions of the sign function, each with its own algebra.

\paragraph{Distributionalism.} Harris and Firth saw that meaning is
distribution, but stopped at statistics because no canonical quotient was
available. The IICA morphism provides the quotient: co-occurrence sets modulo
repetition
and order, with exact set operations. Attention remains the best
\emph{differentiable} approximation of distributional structure; the lattice
is its \emph{exact} image under the quotient.

\paragraph{Category theory.} The bridge is an adjunction: $\Ing$ is a
join-preserving map from the preorder of token contexts (ordered by inclusion)
to the Boolean lattice, and $\Mat$ is its right adjoint on the image, with the
unit of the adjunction being the coverage map and the counit the round-trip
law. IICA is naturality: the diagram commutes independently of the host
language, which is why many paroles can share one langue.

\paragraph{The fifth paper of the series.} Where the first paper
\cite{ewm2026} fixed the principles, the second \cite{companion2026} the
lattice and its worlds, the third \cite{state2026} the trajectory, and the
fourth \cite{attention2026} the map--door--move reading of attention, this
fifth paper states what those four had implicitly built: a linguistic
interpretation in which the lattice is langue, and the whole machine is a
computable structuralism.

% ============================================================================
\section{Conclusion}
\label{sec:conclusion}
% ============================================================================

The paper's claim is that structural linguistics stopped being a philosophical
programme the moment the IICA morphism entered the picture. Attention gave us
parole at scale. The morphism---any deterministic, content-addressed map; a
hash in the reference implementation, chosen for its even distribution---gives
us the quotient from parole to langue. The HLLSet algebra gives langue an
executable form---a Boolean lattice with a
GF(2) ring, whose D/R/N are syntagmatic, whose directed BSS reads succession
and substitution, and whose basis is a system of distinctive oppositions. The
coverage-guarded bridge is
the sign, and IICA is its social contract. The dispatcher is actualization:
the executive side of language, learned.

The circle---statistical linguistics, a content-addressed morphism, algebraic
structures, and back to a fundamental, computable linguistics---is therefore
not an analogy. It is a construction, and it runs.

\section*{Acknowledgments}
This paper was created in collaboration with, and with the assistance of, an
AI system---DeepSeek (deepseek-v4-pro). The five-layer reformulation, the
Saussurean circle, and several of the formulations in this paper emerged in
dialogue with the assistant, whose drafts the authors revised and take full
responsibility for in their final form.

\begin{thebibliography}{99}

\bibitem{saussure1916} F.~de Saussure,
\textit{Cours de linguistique g\'en\'erale}, Payot, 1916.
(Edited posthumously by C.~Bally and A.~Sechehaye.)

\bibitem{hjelmslev1961} L.~Hjelmslev,
\textit{Prolegomena to a Theory of Language}, University of Wisconsin Press,
1961.

\bibitem{harris1954} Z.~S.~Harris,
\textit{Distributional Structure}, Word, 10(2--3):146--162, 1954.

\bibitem{firth1957} J.~R.~Firth,
\textit{A Synopsis of Linguistic Theory, 1930--1955}, Studies in Linguistic
Analysis, 1957.

\bibitem{vaswani2017} A.~Vaswani, N.~Shazeer, N.~Parmar, J.~Uszkoreit,
L.~Jones, A.~N.~Gomez, \L.~Kaiser, and I.~Polosukhin,
\textit{Attention Is All You Need}, NeurIPS, 2017.

\bibitem{ewm2026} A.~Mylnikov, A.~Mylnikov Jr., and D.~Kumar,
\textit{Emerging World Models: A Structural Framework for Self-Organizing
World Representations}, EWM PUB series, 2026.

\bibitem{companion2026} A.~Mylnikov, A.~Mylnikov Jr., and D.~Kumar,
\textit{The EWM Lattice and Its Worlds}, EWM PUB series, 2026.

\bibitem{state2026} A.~Mylnikov, A.~Mylnikov Jr., and D.~Kumar,
\textit{Transformations and Trajectories of a Content-Addressed State
Machine}, EWM PUB series, 2026.

\bibitem{attention2026} A.~Mylnikov, A.~Mylnikov Jr., and D.~Kumar,
\textit{Attention on Lattice Traversal Is All You Need}, EWM PUB series, 2026.

\bibitem{birkhoff1940} G.~Birkhoff,
\textit{Lattice Theory}, American Mathematical Society, 1940.

\bibitem{maclane1971} S.~Mac Lane,
\textit{Categories for the Working Mathematician}, Springer, 1971.

\end{thebibliography}

\end{document}
